\section{Theoretical error analysis of SR versus RN}
\label{sec:analysis}

The GPT-2 architecture contains projections with very different reduction
lengths---the MLP down-projection accumulates over $n=3072$ terms, whereas the
attention projections accumulate over $n=768$---and the language-model head
applies the softmax directly to the perturbed logits.  These structural
differences create distinct opportunities for SR's variance--drift trade-off.
This section analyzes both regimes:
first, establishing a probabilistic forward error bound for the affine
projection layers (\code{attn\_c\_proj} and \code{mlp\_c\_proj}) via a vector Bienaym\'e--Chebyshev
construction (\cref{sec:analysis-mlp}); second, deriving an analytical decomposition
of cross-entropy loss change into drift and variance components at the output softmax
(\cref{sec:analysis-softmax}); and third, empirically validating this trade-off on DistilGPT-2
(\cref{sec:analysis-empirical-drift}).
Throughout, statements concern the precision at which a reduction \emph{accumulates}.
On physical hardware, a kernel nominally in \code{bfloat16} may store in \code{bfloat16},
multiply into a wider product, and accumulate in \code{binary32}; under PRISM,
every elementary operation rounds at the same virtual precision, $u_{\rm sto} = u_{\rm mul} = u_{\rm acc} = 2^{1-t}$,
so the relative ULP scale $u$ below is $2^{1-t}$ for the significand precision $t$
swept by our harness.

\subsection{Forward error analysis of linear projection layers}
\label{sec:analysis-mlp}

The multilayer perceptron (MLP) block in the GPT-2 architecture consists of two
consecutive linear projections surrounding an element-wise activation function.
Following the row-vector convention of \cref{eq:mlp},
\begin{equation*}
  \text{MLP}(x) = \phi(x W_1 + b_1) W_2 + b_2,
\end{equation*}
where $\phi$ is the GELU activation function~\citep{hendrycks2016gelu},
$W_1 \in \mathbb{R}^{d_{\rm model} \times d_{\rm ff}}$ is the expansion matrix (\code{mlp\_c\_fc}), and
$W_2 \in \mathbb{R}^{d_{\rm ff} \times d_{\rm model}}$ is the projection matrix (\code{mlp\_c\_proj}).
In DistilGPT-2, $d_{\rm model} = 768$ and $d_{\rm ff} = 3072$.

\paragraph{Limitations of backward error frameworks for nonlinear activations.}
The backward error framework of \citet{beuzeville2025deterministic}
analyzes multi-layer networks by absorbing rounding errors into equivalent weight
perturbations $\Delta W$. For a layer followed by an activation $\phi$, the
componentwise backward error bound includes a term proportional to
$l_\phi / \kappa_\phi(z)$, where $\kappa_\phi(z) = |z \phi'(z) / \phi(z)|$ is the
relative condition number of $\phi$ and $l_\phi u$ bounds the evaluation error
of $\phi$.

For GELU, this framework does not yield a uniform backward guarantee due to the
vanishing derivative at its local minimum. Specifically, at $z^* \approx -0.752$,
$\phi'(z^*) = 0$ while $\phi(z^*) \approx -0.170 \neq 0$. Consequently, the
condition number vanishes ($\kappa_\phi(z^*) = 0$), causing the backward error bound
$l_\phi / \kappa_\phi(z)$ to diverge in the neighborhood of $z^*$.

Although each ideal SR step is conditionally unbiased, propagating SR perturbations
through a nonlinear activation generally introduces composition bias
\citep[Section~4]{elarar2024bounds}. To isolate the reduction error from this
nonlinear composition, we focus on the affine down-projection \code{mlp\_c\_proj} (and, by extension, the attention projection \code{attn\_c\_proj}):
\begin{equation*}
  y = h W + b,
\end{equation*}
where $h \in \mathbb{R}^n$ is the input row vector, $W \in \mathbb{R}^{n \times m}$ is the projection matrix,
$b \in \mathbb{R}^m$ is the bias vector, and $y \in \mathbb{R}^m$ is the output ($m = d_{\rm model} = 768$).
For \code{mlp\_c\_proj}, $n = d_{\rm ff} = 3072$ and $h = \phi(x W_1 + b_1)$; for \code{attn\_c\_proj}, $n = d_{\rm model} = 768$ and $h$ is the concatenated multi-head attention output.
Because these projection layers have no subsequent activation ($\phi = \mathrm{id}$), the forward
error $e = \hat{y} - y$ directly quantifies the perturbation injected into the
residual stream.

\paragraph{Error model.}
Each component of $y$ is evaluated as a left-to-right sequential inner product
followed by a bias addition:
\begin{equation*}
  y_i = \sum_{j=1}^{n} h_j \, W_{ji} + b_i, \qquad i = 1, \dots, m,
\end{equation*}
totaling $2n$ operations per output component with a longest error-product chain
of $n + 1$ rounding factors (one multiplication, $n - 1$ accumulations, and the
final bias addition).

Let $u = 2^{1-t}$ denote the relative ULP scale at virtual significand precision $t$.
Assuming no overflow, underflow, or NaNs occur, every elementary operation satisfies
$\fl(a \mathbin{\mathrm{op}} b) = (a \mathbin{\mathrm{op}} b)(1 + \delta)$, where:
\begin{itemize}[leftmargin=1.5em]
  \item under RN, $|\delta| \le u/2$;
  \item under SR (mode-1, nearness), $|\delta| \le u$, with the mean-independence
        property $\mathbb{E}[\delta_k \mid \delta_1, \dots, \delta_{k-1}] = 0$~\citep[Lemma~5.2]{connolly2021stochastic}.
\end{itemize}

\paragraph{Deterministic forward error bound (RN).}
Under deterministic RN, rounding errors may align in the worst case. Applying Higham's
exact forward error bound for inner products~\citep[Lemma~3.1]{higham2002accuracy}, each component
and then the vector satisfy
\begin{align*}
  |\hat{y}_i - y_i| &\le \left( (1 + u/2)^{n+1} - 1 \right) \left( \sum_{j=1}^{n} |h_j| |W_{ji}| + |b_i| \right), \\
  \|\hat{y} - y\|_2 &\le \left( (1 + u/2)^{n+1} - 1 \right) \, \| |h| \, |W| + |b| \|_2.
\end{align*}
In the high-precision regime where $nu \ll 1$, the coefficient linearizes to $\approx \frac{n+1}{2} u$;
when $nu \gg 1$, it grows exponentially, reaching $\approx 8.11 \times 10^{80}$ at $t = 4$ for $n=3072$.

\paragraph{Probabilistic rounding and scalar bounds.}
In the first-order regime, mean-independence allows SR forward error bounds
proportional to $\mathcal{O}(\sqrt{n} \, u)$ rather than the $\mathcal{O}(n \, u)$
deterministic RN bound. Existing probabilistic frameworks for bounding scalar
inner product errors $e_i = \hat{y}_i - y_i$ follow one of two strategies:
\begin{itemize}[leftmargin=1.5em]
  \item \emph{Azuma--Hoeffding (AH) method}~\citep{connolly2021stochastic,ipsen2020probabilistic,de2025error}:
        Martingale bounds provide exponential probability tails, which tolerate a union bound
        over dimensions with only a mild $\mathcal{O}(\sqrt{\ln m})$ penalty. However, the bound on
        the increments contains a factor $(1 + u)^{n+1}$~\citep[Section~7]{elarar2024bounds}.
        When $nu \gg 1$, this factor explodes exponentially,
        rendering AH uninformative at low precisions.
  \item \emph{Bienaym\'e--Chebyshev (BC) method}~\citep{elarar2022error,elarar2024bounds}:
        Bounding the variance of the accumulated error avoids this exponential explosion:
        because the precision enters squared in the variance bound,
        $(1 + u^2)^{n+1} - 1 \approx (n+1) u^2$ remains small even at low precisions where $nu \gg 1$.
\end{itemize}
The core variance bound for error products under SR is given by the following lemma.

\begin{lemma}[{\citealp[Lemma~3.1]{elarar2022error}}]
  \label[lemma]{lem:sr-variance}
  Let $\varphi = \prod_{k=1}^{K}(1 + \delta_k)$, where $\delta_k$ are SR-nearness rounding
  errors satisfying $|\delta_k| \le u$ and $\mathbb{E}[\delta_k \mid \delta_1, \dots, \delta_{k-1}] = 0$.
  Then $\mathbb{E}[\varphi] = 1$ and $\Var(\varphi) \le (1 + u^2)^K - 1$.
\end{lemma}

Applying Chebyshev's inequality to each scalar component $e_i$ yields an error bound scaling with $\frac{1}{\sqrt{\alpha}}\sqrt{(1 + u^2)^{n+1} - 1}$.

\paragraph{Vector $L_2$-norm bound: combining BC with a multivariate Markov bound.}
Extending the scalar analysis to the full output vector $e \in \mathbb{R}^m$ exposes a limitation in existing bounds: AH survives the union bound across $m$ components but fails at low precision $t$, whereas scalar BC remains stable at low precision but grows rapidly under the union bound. Because Chebyshev's inequality has a polynomial tail, setting the per-component failure probability to $\alpha / m$ inflates the confidence multiplier by $\sqrt{m}$, from $1/\sqrt{\alpha} \approx 4.47$ to $\sqrt{m/\alpha} \approx 124$ ($m = 768$, $\alpha = 0.05$). This penalty inflates the bound to $\approx 6870 \, u$, wiping out SR's advantage over RN.

To circumvent this limitation, our approach combines the low-precision robustness of BC with a single multivariate Markov step applied directly to the expected squared $L_2$-norm $\mathbb{E}[\|e\|_2^2]$. This bypasses the union bound, keeping the confidence multiplier $1/\sqrt{\alpha}$ independent of the output dimension $m$.

\begin{theorem}[Vector Bienaym\'e--Chebyshev Bound]
  \label{thm:vector-bc}
  Under SR-nearness, the forward error $e = \hat{y} - y$ of the projection $y = h W + b$
  satisfies, for any $\alpha \in (0, 1)$, with probability at least $1 - \alpha$:
  \begin{equation*}
    \|\hat{y} - y\|_2 \le \frac{1}{\sqrt{\alpha}} \sqrt{(1 + u^2)^{n+1} - 1} \, \| |h| \, |W| + |b| \|_2.
  \end{equation*}
  The confidence multiplier $1/\sqrt{\alpha}$ is independent of the output dimension $m$.
\end{theorem}

\begin{proof}
  The error in output component $i$ is
  \begin{equation*}
    e_i = \sum_{j=1}^{n} h_j W_{ji} (\varphi_{ij} - 1) + b_i (\varphi_{b,i} - 1),
  \end{equation*}
  where $\varphi_{ij}$ and $\varphi_{b,i}$ represent the error products along the accumulation
  chains, containing at most $n + 1$ and $1$ rounding factors, respectively. By mean-independence,
  $\mathbb{E}[\varphi_{ij}] = 1$, $\mathbb{E}[\varphi_{b,i}] = 1$, and $\mathbb{E}[e_i] = 0$.

  Although the terms $\varphi_{ij} - 1$ share common rounding errors across additions,
  Minkowski's inequality in $L_2$ bounds the standard deviation of their linear combination:
  \begin{equation*}
    \sqrt{\Var(e_i)} \le \sum_{j=1}^{n} |h_j| |W_{ji}| \sqrt{\Var(\varphi_{ij})} + |b_i| \sqrt{\Var(\varphi_{b,i})}.
  \end{equation*}
  Applying \cref{lem:sr-variance} with $\Var(\varphi_{ij}) \le (1 + u^2)^{n+1} - 1$ and
  $\Var(\varphi_{b,i}) \le u^2 \le (1 + u^2)^{n+1} - 1$ yields
  \begin{equation*}
    \Var(e_i) \le \left( (1 + u^2)^{n+1} - 1 \right) \left( \sum_{j=1}^{n} |h_j| |W_{ji}| + |b_i| \right)^2.
  \end{equation*}
  By linearity of expectation, the expected squared $L_2$-norm satisfies
  \begin{equation*}
    \mathbb{E}[\|e\|_2^2] = \sum_{i=1}^{m} \Var(e_i) \le \left( (1 + u^2)^{n+1} - 1 \right) \, \| |h| \, |W| + |b| \|_2^2.
  \end{equation*}
  Applying Markov's inequality to the squared $L_2$-norm $\|e\|_2^2$ (which constitutes the
  multivariate Bienaym\'e--Chebyshev inequality) gives
  $\mathbb{P}\left( \|e\|_2^2 \ge \frac{1}{\alpha} \mathbb{E}[\|e\|_2^2] \right) \le \alpha$.
  Taking the square root completes the proof.
\end{proof}

\paragraph{Numerical evaluation on DistilGPT-2.}
While \cref{thm:vector-bc} holds for any affine projection of reduction length $n$,
we evaluate the bounds concretely on the feedforward down-projection \code{mlp\_c\_proj}
($n = d_{\rm ff} = 3072$, $m = d_{\rm model} = 768$) at 95\%
confidence ($\alpha = 0.05$, $1/\sqrt{\alpha} = \sqrt{20} \approx 4.47$).
Because \code{mlp\_c\_proj} has four times the reduction length of \code{attn\_c\_proj} ($n = 768$),
it has the longest accumulation chains and therefore tests the error bounds in the most demanding regime.
In the limiting first-order regime $u\to0$ (equivalently, $t\to\infty$):
\begin{itemize}[leftmargin=1.5em]
  \item RN deterministic: $(1 + u/2)^{n+1} - 1 \approx \frac{n+1}{2} u \approx 1537 \, u$;
  \item SR vector BC: $\frac{1}{\sqrt{\alpha}} \sqrt{(1 + u^2)^{n+1} - 1} \approx \sqrt{20} \sqrt{n+1} \, u \approx 248 \, u$.
\end{itemize}
The SR coefficient is asymptotically smaller than the RN coefficient by a factor
of approximately $6.2$. At finite precision the exact ratio approaches this
limit from above: it is $14.37$ at $t=11$ and $6.35$ at $t=16$.

\begin{figure}[t]
  \centering
  \safeincludegraphics[width=0.75\linewidth]{figures/rn_vs_sr_bounds.pdf}
  \caption{Deterministic RN and SR forward-error bound
  coefficients for \code{mlp\_c\_proj} in DistilGPT-2 ($n=3072$).}
  \label{fig:validity-bounds}
\end{figure}

As shown in \cref{fig:validity-bounds}, the deterministic RN bound suffers from exponential error
compounding at low precision ($nu \gg 1$), reaching $2.43 \times 10^{10}$ at $t = 7$ and $8.11 \times 10^{80}$ at $t = 4$.
In contrast, because precision enters squared under stochastic rounding, the SR
vector BC bound grows much more slowly ($4.73$ at $t = 7$). At $t=11$ the exact RN and SR
coefficients are $3.482$ and $0.242$, a ratio of $14.37$; the ratio decreases
toward $6.2$ only in the limit $t\to\infty$.

For fixed $\alpha$ in the first-order regime, \cref{thm:vector-bc} yields an
$\mathcal{O}(\sqrt{n}u)$ bound for SR, whereas deterministic RN scales as
$\mathcal{O}(nu)$. Both bounds are worst-case over the inputs and normalize by
$\| |h| \, |W| + |b| \|_2$, but they treat arithmetic errors
differently. Under RN, rounding errors are deterministic functions of the data
that can align constructively in the worst case. Under SR, mean-independence
yields the $\sqrt{n}$ factor via the Bienaym\'e--Chebyshev variance bound. While
both results are upper bounds, and realized RN errors rarely achieve worst-case
alignment, the comparison illustrates the fundamental contrast between the two modes:
deterministic rounding permits systematic linear drift along the reduction chain,
whereas stochastic rounding restricts error growth to a sub-linear
$\mathcal{O}(\sqrt{n}u)$ envelope.

The downstream captures in \cref{sec:analysis-empirical-drift} show larger RN than SR
sample logit drift after reducing \code{mlp\_c\_proj}. These propagated
measurements do not directly test the single-reduction bounds above.
